\begin{longtable}{@{}l p{12.5cm}@{}}
  \caption{Skills per-instance answer comparison ($\Delta \leq -75\%$) for the 1 most divergent items: same question, each model's answer and reasoning.}\label{tab:cmp_skills_neg75} \\
  \toprule
  \endfirsthead
  \toprule
  \endhead
  \multicolumn{2}{@{}l}{\textbf{83180} \quad $\Delta = -80\%$} \\
  \textbf{Question} & You are evaluating a learner's responses on daily living skills training for a diagnostic assessment. Based on the learner's incorrect answers and their written explanations, analyze and classify the most likely underlying cause(s) of their errors (e.g., lack of prerequisite knowledge, misunderstanding of procedure, attention/memory lapse, or affective resistance). Describe specifically which prerequisite concepts appear missing, which steps or ideas are misunderstood, and how the learner's affective responses (motivation, anxiety, or refusal) may have contributed. Finally, summarize recommended next instructional focus areas to address the diagnosed causes and support the learner's participation in future lessons. \\
  \textbf{Reference} & After reviewing the learner's incorrect responses and written explanations, the primary causes of error appear to be a combination of missing prerequisite knowledge and misunderstandings of procedural steps, with a secondary contribution from mild anxiety that reduced engagement. Specifically: \newline 1) Missing prerequisite knowledge: The learner consistently chose actions that indicate they do not fully grasp the concept of sequencing in multi-step tasks. For example, when asked to describe how to prepare a simple snack, they listed steps out of logical order (putting a plate away before assembling ingredients), suggesting they lack an internalized routine schema for task sequencing and the purpose of preparatory steps. They also showed limited vocabulary for safety-related terms (e.g., confused hot with warm), indicating gaps in basic safety knowledge needed for independent daily living tasks. \newline 2) Misunderstanding of procedure: In responses about personal hygiene, the learner described rinsing with a large splash rather than a controlled washing motion and omitted critical substeps such as drying hands and turning off the tap. This pattern shows a surface-level awareness of the goal (hands should be clean) but misunderstanding of the standard procedure and substep importance. Their written explanation emphasized doing the action quickly rather than following each step, which points to an incomplete mental model of why each substep matters. \newline 3) Affective factors: The learner's language in explanations included phrases like I don't like doing that and I get nervous if it takes too long, indicating low motivation and task-related anxiety. These affective responses likely contributed to rushed answers and omission of steps. There was no evidence of outright refusal, but the anxious wording suggests that affect decreased their persistence and attention to detail. \newline Recommended instructional focus areas: \newline - Teach and practice sequencing explicitly: use visual task analysis (picture or written checklists) and guided chaining (forward or backward chaining) to build a stable routine for multi-step activities. Begin with very short chains (2--3 steps) and gradually increase complexity as the learner demonstrates mastery. \newline \textit{[\ldots\ truncated \ldots]} \newline - Address affective barriers: incorporate brief coping strategies for anxiety (simple deep-breathing before tasks), positive reinforcement for completing all steps, and offer choices within tasks to increase motivation (e.g., choose snack components). Start with high-success tasks to build confidence. \newline Summary diagnostic statement: The learner's errors stem mainly from gaps in sequencing knowledge and incomplete procedural understanding, exacerbated by low motivation and mild anxiety that lead to rushed or omitted steps. Instruction should prioritize explicit sequencing instruction, targeted safety vocabulary, scaffolded procedural practice, and brief affective supports to improve participation and reduce errors in future lessons. \\
  \midrule
  \textbf{Kimi-K2.6} & score 10/10 \\
  & [9-10] The student's answer is highly systematic and detailed, providing accurate classifications with specific rationale and examples for each category (prerequisites, procedural misunderstandings, affective factors). The instructional focus areas are clear, actionable, and directly address the diagnosed causes, meeting all criteria for the highest band. \\
  \textbf{Innospark-1T} & score 2/10 \\
  & Let me think through this carefully. \newline The question/task: You are evaluating a learner's responses on daily living skills training for a diagnostic assessment. Based on the learner's incorrect answers and their written explanations, analyze and classify the most likely underlying cause(s) of their errors (e.g., lack of prerequisite knowledge, misunderstanding of procedure, attention/memory lapse, or affective resistance). Descr \newline My reasoning: Diagnostic Assessment Analysis Framework \newline Since no specific learner responses were provided, I'll present a comprehensive framework for analyzing daily living skills errors, with illustrative examples. General Diagnostic Framework \newline Step 1: Categorize Error Types \newline | Error Category | Indicators in Written Explanations | \newline \textit{[\ldots\ truncated \ldots]} \newline ``` \newline To provide a specific analysis, please share the learner's actual responses and written explanations. \\
  \specialrule{1.2pt}{8pt}{8pt}
  \bottomrule
\end{longtable}